\documentclass[10pt,a4paper]{article}
\usepackage{amsmath,amssymb,amsthm,graphicx}
\usepackage{amsmath}
\usepackage{amsfonts}
\usepackage{amssymb}
\usepackage{amsthm}
\usepackage{enumerate}
\newtheorem{theorem}{Theorem}[section]
\newtheorem{conjecture}{Conjecture}[section]
\newtheorem{lemma}[theorem]{Lemma}
%\theoremstyle{definition}
\newtheorem{definition}[theorem]{Definition}
\newtheorem{corollary}[theorem]{Corollary}
\newtheorem{example}[theorem]{Example}

\newcommand{\summ}{\displaystyle\sum}
\newcommand{\prodd}{\displaystyle\prod}
\newcommand{\F}{\mathbb{F}}
\newcommand{\E}{\mathbf{E}}
\newcommand{\calF}{\mathcal{F}}
\newcommand{\Feps}{{\calF}_{\epsilon}}
\newcommand{\R}{\mathbb{R}}
\newcommand{\mon}{\mathcal{M}}
\newcommand{\supp}{\mathsf{supp}}
\newcommand{\wt}{\mathsf{wt}}
\newcommand{\poly}{\mathsf{poly}}
\newcommand{\one}{{\mathbf{1}}}
\newcommand{\qbin}[3]{\begin{bmatrix}{#1}\\ {#2}\end{bmatrix}_{#3}}
\title{A nontrivial degree upper bound for square-root polynomials for every odd prime}
\author{Kiran Kedlaya, Swastik Kopparty}
\date{doi:10.4230/LIPIcs.CCC.2024.25; CC BY 4.0}
\begin{document}
\maketitle
\setcounter{section}{5}
\noindent\textit{Allowed standard prerequisites (editorial).}\begin{itemize}
\item Polynomial vanishing sets over finite fields and formal differentiation. Author Sections 1-3; original full bibliography retained
\item Weil bound for mixed character sums and Fourier L1 bounds for intervals. Author Section 5.1 and original Schmidtbook reference; local equidistribution proof included
\end{itemize}\medskip
\begin{theorem}
 For all odd primes $p$, for $t = o\left(\frac{\sqrt{p}}{\log^2 p} \right)$, there is a polynomial of degree $\frac{p}{2} - t$ which computes square roots.
\end{theorem}

Let $p$ be an odd prime, and let $\F_p$ be the finite field with $p$ elements.

We say $f(X) \in \F_p[X]$ {\em computes square roots} in $\F_p$, if
for all nonzero perfect squares $a \in \F_p$, we have:
$$f(a)^2 = a.$$ In other words, for each nonzero perfect square $a \in \F_p$, $f(a)$ is one of its two square roots.

\section{Upper bounds for general $p$}


In this section, we give a slightly nontrivial upper bound on the degree of polynomials computing square roots for all $p$.
We will show that there is a polynomial with degree somewhat less than $\frac{p}{2}$ which computes square roots.

\begin{proof}
Let $m = (p-1)/2$. Let $S \subseteq \F_p$ be the set of nonzero perfect squares,
and note that $|S| = m$.
For each $\alpha \in S$, let $\delta_\alpha(X) \in \F_p[X]$ be
the unique polynomial of degree $\leq (m-1)$ such that for all $\beta\in S$:
$$ \delta_\alpha(\beta) =  \begin{cases} 1  & \beta = \alpha, \\ 0 & \beta \neq \alpha. \end{cases}$$
Explicitly, we have:
$$\delta_\alpha(X) = \frac{1}{m} \left( \left(\frac{X}{\alpha}\right)^{m-1} + \left(\frac{X}{\alpha}\right)^{m-2} + \ldots + \frac{X}{\alpha} + 1 \right).$$


Then given a function $u: S \to \F_p$, 
the unique polynomial $f(X) \in \F_p[X]$ of degree at most $m-1$
such that $f(\alpha) = u(\alpha)$ for all $\alpha \in S$ is given by:
$$f(X) =  \sum_{\alpha \in S} u(\alpha) \delta_\alpha(X).$$

Our goal is to pick $u$ where each $u(\alpha)$ is one of the two square roots of $\alpha$ so that many of the leading coefficients of $f(X)$ equal $0$.

We now use the structure of $S$. Let $g$ be a generator of $\F_p^*$. 
Then $S = \{g^{2j} \mid  0 \leq j < (p-1)/2 \}$. Furthermore, 
for $\alpha = g^{2j} \in S$, one of the two square roots of $\alpha$
is $g^{j}$.

Thus, our problem can be reformulated as choosing a function
$v: S \to \{\pm 1 \}$ such that:
$$ f(X) = \sum_{j=0}^{(p-1)/2} v\left(g^{2j}\right)\cdot g^j \cdot \delta_{g^{2j}}(X)$$
has many leading coefficients equal to $0$.

Observe that the coefficient of $X^{m-i}$ in $f(X)$ equals:
$$ \frac{1}{m} \sum_{j=0}^{(p-1)/2} v(g^{2j}) g^{j} \left( \frac{1}{g^{2j}} \right)^{m-i} = 
 \frac{1}{m} \sum_{i=0}^{(p-1)/2} v(g^{2j}) g^{(2i+1)j}.$$


 Thus, to get a polynomial $f(X)$ of degree $< m-t$, we want to find a vector $v \in \{\pm 1\}^m$ that lies in the kernel of the Vandermonde-type matrix $M \in \F_p^{t \times m}$, where:
 $$ M_{i,j} = g^{(2i+1)j}.$$
 (The row index $i$ runs from $1$ to $t$, the column index $j$ runs from $0$ to $m-1$.)
 
 For later use, for a vector $y \in \F_p^t$, we define $P_y(Z) \in \F_p[Z]$ to be the polynomial:
 $$ P_y(Z) = \sum_{i = 1}^{t} y_i Z^{2i+1}.$$
 Thus for $j \in \{0, 1, \ldots, m-1 \}$, the $j$th entry of $M^Ty$ equals $P_y(g^{j})$.


 To show that there exists the desired $\pm 1$ vector, we count the number of such vectors using Fourier analysis. Let $\omega$ be a $p$th root of unity in $\mathbb C$. The number of such $\pm 1$ vectors equals:
 \begin{align*}
 N &= \sum_{v \in \{\pm 1\}^m } \one_{M v = 0}\\
 &= \sum_{v \in \{\pm 1\}^m }  \mathbb E_{y \in \F_p^t}  \left[ \omega^{\langle y, Mv \rangle} \right]  \\
 &= \mathbb E_y \left[ \sum_v  \omega^{\langle M^T y, v \rangle}   \right]\\
 &= \mathbb E_y \left[ \sum_v  \omega^{\sum_{j=0}^{m-1} (M^T y)_j \cdot  v_j}   \right]\\
 &= \mathbb E_y \left[ \sum_v   \prod_{j=0}^{m-1} \omega^{(M^T y)_j \cdot  v_j}   \right]\\
 &= \mathbb E_y \left[ \sum_v \prod_{j=0}^{m-1} \omega^{P_y(g^{j}) \cdot  v_j} \right].
 \end{align*}
 
%The number of ways of getting the top $t$ coefficients to equal $0 \in \F_p^t$ equals:
%$$\mathbb E_{y \in \F_p^t} \left[\sum_{u \in \{\pm 1\}^m} \omega^{\sum_{\alpha} u(\alpha) Y(\alpha) }\right].$$ 
%Here $Y(X) \in \F_p[X]$ is the polynomial of degree $t-1$ with coefficients $y$.

For $y = 0$, the expression inside the expectation equals $2^m$.
We will show that for the remaining $p^t -1$ values of $y$,
the expression inside the expectation is very small. 

Fix any $y \neq 0$. 
The expression inside the expectation equals:
\begin{align}
 \label{weilval}
\sum_{v \in \{\pm 1\}^m} \prod_{j=0}^{m-1} \omega^{P_y(g^{j}) \cdot  v_j} = \prod_{j=1}^m \left(\omega^{P_y(g^{j})} + \omega^{-P_y(g^{j})} \right).
\end{align}

The next lemma (which uses the Weil bounds on mixed character sums) shows that for any nonzero $y$, the evaluations of the polynomial $P_y$ at $\{1, g, g^2, \ldots, g^{m-1}\}$ are well distributed in $\F_p$.
\begin{lemma}
\label{lem:equidist}
Let $0 \leq \alpha <  \beta \leq 1$.
 Let $y \in \F_p^t \setminus \{0\}$. Then:
 $$ \Pr_{j \in \{0,1, \ldots, m-1\}}\left[ P_y(g^j) \in [\alpha p, \beta p] \right] = (\beta - \alpha )  + O\left(\frac{t\log^2 p}{\sqrt{p}} \right).$$
\end{lemma}

Assuming the lemma, we get that for $t = o\left(\frac{\sqrt{p}}{\log^2 p}\right)$, the product in Equation \eqref{weilval} is at most $2^m \cdot \exp(-m)$.
Thus:
\begin{align*} N &\geq \frac{2^m}{p^t}  - \max_{y \neq 0} \left| \prod_{j=1}^m \left(\omega^{P_y(g^{j})} + \omega^{-P_y(g^{j})} \right) \right|\\
 &\geq \frac{2^m}{p^t}  - O(2^{m} \exp(-m))\\
 &\geq 2^m \left(\frac{1}{p^t} - \exp(-m)  \right)\\
 & > 0,
\end{align*}
where the penultimate inequality holds since
$$t = o\left( \frac{\sqrt{p}}{\log^2 p} \right) \ll \frac{p}{\log p} = \Theta\left( \frac{m}{\log p} \right).$$

This completes the proof.
\end{proof}


\subsection{Distribution of values of $P_y(g^j)$}

We now prove Lemma~\ref{lem:equidist}.

\begin{proof}
 Let $I$ be the interval $[\alpha p, \beta p]$. Let $J$ be the set
 $\{1, g, g^2, \ldots, g^{m-1}\}$. Then the probability in the statement of the lemma equals:
 \begin{align}
 \label{pval}
 \frac{1}{m} \sum_{ z \in \F_p}  {\one}_{I}(P_y(z)){\one}_J(z).
 \end{align}
 
 $I$ is an interval in the additive group of $\F_p$
 By standard results, if we expand $\one_I$ in its additive Fourier expansion:
 $$ \one_I = \sum_\psi \widehat{\one}_I(\psi) \psi$$
 (where the $\psi$ are the additive characters), then:
 \begin{align}
 \label{l1bound1}
 \sum_{\psi } |\widehat{\one}_I(\psi)| \leq O( \log p).
 \end{align}
 
 Similarly, $J$ is an interval in the multiplicative group of $\F_p$.
 If we expand $\one_J$ in its additive Fourier expansion:
 $$ \one_J = \sum_\psi \widetilde{\one}_J(\chi) \chi$$
 (where the $\chi$ are the multiplicative characters), then:
 \begin{align}
 \label{l1bound2}
 \sum_{\chi} |\widetilde{\one}_I(\chi)| \leq O( \log p).
 \end{align}
 
 Then the probability in Equation~\eqref{pval} equals:
 \begin{align*} &\frac{1}{m}\sum_{z} \left( \sum_{\psi} \widehat{\one}_I(\psi) \psi(P_y(z)) \right) \left( \sum_{\chi} \widetilde{\one}_J(\chi) \chi(z) \right)\\
  &= \frac{1}{m}\sum_{\psi, \chi} \widehat{\one}_I(\psi) \widetilde{\one}_J(\chi) \left(\sum_{z} \psi(P_y(z)) \chi(z) \right)\\
  &= \frac{1}{m} \left( \frac{|I|}{p}\cdot \frac{|J|}{p} \cdot p + O\left(\sum_{(\psi, \chi) \neq (1,1)} |\widehat{\one}_I(\psi)|\cdot |\widetilde{\one}_J(\chi)| \cdot \left|\sum_{z} \psi(P_y(z)) \chi(z) \right|\right) \right)\\
  &= (\beta-\alpha) + \frac{1}{m} \cdot  O\left(\sum_{(\psi, \chi) \neq (1,1)} |\widehat{\one}_I(\psi)| \cdot | \widetilde{\one}_J(\chi) | \cdot \left|\sum_{z} \psi(P_y(z)) \chi(z) \right|\right) \\
 \end{align*}
The Weil bound for mixed character sums (see~\cite{Schmidtbook}) shows that whenever
$(\psi, \chi)$ are not both trivial characters, the inner expression is bounded:
$$\left|\sum_{z} \psi(P_y(z)) \chi(z) \right| \leq O( t \sqrt{p}).$$
Combined with the bounds in Equations~\eqref{l1bound1}, \eqref{l1bound2}, we get the desired bound on the probability.
\end{proof}

\begin{thebibliography}{VZGG13}

\bibitem[ADN03]{ADN-sqrt}
Simon~Joseph Agou, Marc Del{\'e}glise, and Jean-Louis Nicolas.
\newblock Short polynomial representations for square roots modulo p.
\newblock {\em Designs, Codes and Cryptography}, 28(1):33--44, 2003.

\bibitem[AS97]{ASlowdegree}
Sanjeev Arora and Madhu Sudan.
\newblock Improved low-degree testing and its applications.
\newblock In {\em Proceedings of the twenty-ninth annual ACM symposium on
  Theory of computing}, pages 485--495, 1997.

\bibitem[Bir00]{Biro}
Andr{\'a}s Bir{\'o}.
\newblock On polynomials over prime fields taking only two values on the
  multiplicative group.
\newblock {\em Finite Fields and Their Applications}, 6(4):302--308, 2000.

\bibitem[BS96]{BachShallit}
Eric Bach and Jeffrey~Outlaw Shallit.
\newblock {\em Algorithmic number theory: Efficient algorithms}, volume~1.
\newblock MIT press, 1996.

\bibitem[Car11]{Carella-sqrt}
NA~Carella.
\newblock Formulas for the square root modulo p.
\newblock {\em arXiv preprint arXiv:1101.4605}, 2011.

\bibitem[CKL15]{CKL-nthroot}
Seunghwan Chang, Bihtnara Kim, and Hyang-Sook Lee.
\newblock Polynomial representations for n-th roots in finite fields.
\newblock {\em Journal of the Korean Mathematical Society}, 52(1):209--224,
  2015.

\bibitem[CS00]{CS-dlog}
Don Coppersmith and Igor Shparlinski.
\newblock On polynomial approximation of the discrete logarithm and the
  diffie—hellman mapping.
\newblock {\em Journal of Cryptology}, 13:339--360, 2000.

\bibitem[GR05]{GR-listrecovery}
Venkatesan Guruswami and Atri Rudra.
\newblock Limits to list decoding reed-solomon codes.
\newblock In {\em Proceedings of the thirty-seventh annual ACM symposium on
  Theory of computing}, pages 602--609, 2005.

\bibitem[GS98]{RSlist}
Venkatesan Guruswami and Madhu Sudan.
\newblock Improved decoding of reed-solomon and algebraic-geometric codes.
\newblock In {\em Proceedings 39th Annual Symposium on Foundations of Computer
  Science (Cat. No. 98CB36280)}, pages 28--37. IEEE, 1998.

\bibitem[Kal91]{KaltofenHIT}
Erich Kaltofen.
\newblock Effective noether irreducibility forms and applications.
\newblock In {\em Proceedings of the twenty-third annual ACM symposium on
  Theory of Computing}, pages 54--63, 1991.

\bibitem[Mas84]{abcthm2}
Richard~Clive Mason.
\newblock {\em Diophantine equations over function fields}, volume~96.
\newblock Cambridge University Press, 1984.

\bibitem[MW86]{MW-dlog}
Gary Mullen and David White.
\newblock A polynomial representation for logarithms in gf (q).
\newblock {\em Acta arithmetica}, 3(47):255--261, 1986.

\bibitem[Sch06]{Schmidtbook}
Wolfgang~M Schmidt.
\newblock {\em Equations over finite fields: an elementary approach}, volume
  536.
\newblock Springer, 2006.

\bibitem[Sha73]{ShanksTonelli1}
Daniel Shanks.
\newblock Five number-theoretic algorithms.
\newblock In {\em Proceedings of the Second Manitoba Conference on Numerical
  Mathematics (Winnipeg), 1973}, 1973.

\bibitem[Sho09]{Shoup}
Victor Shoup.
\newblock {\em A computational introduction to number theory and algebra}.
\newblock Cambridge university press, 2009.

\bibitem[Sto81]{abcthm1}
W~Wilson Stothers.
\newblock Polynomial identities and hauptmoduln.
\newblock {\em The Quarterly Journal of Mathematics}, 32(3):349--370, 1981.

\bibitem[Ton91]{ShanksTonelli2}
Alberto Tonelli.
\newblock Bemerkung {\"u}ber die aufl{\"o}sung quadratischer congruenzen.
\newblock {\em Nachrichten von der K{\"o}nigl. Gesellschaft der Wissenschaften
  und der Georg-Augusts-Universit{\"a}t zu G{\"o}ttingen}, 1891:344--346, 1891.

\bibitem[VZGG13]{GerhardGathen}
Joachim Von Zur~Gathen and J{\"u}rgen Gerhard.
\newblock {\em Modern computer algebra}.
\newblock Cambridge university press, 2013.

\bibitem[WB83]{BerlekampWelch}
Lloyd~R Welch and Elwyn~R Berlekamp.
\newblock Error correction for algebraic block codes.
\newblock {\em US patent}, (4,633,470), 1983.

\bibitem[Win02]{W-dlog}
Arne Winterhof.
\newblock Polynomial interpolation of the discrete logarithm.
\newblock {\em Designs, Codes and Cryptography}, 25(1):63--72, 2002.

\end{thebibliography}
\end{document}
